\section{Continuity of one resource}\label{sec:continuity}

\subsection{Contracts, reads and captures}
\begin{definition}[Contract]\label{def:contract}
Fix a resource, its request context (including content negotiation) and the
authority whose state a read is meant to describe. A \emph{contract} is a
finite directed graph $G=(V,E)$ of every server state and transition permitted
in that scope, including those of every writer and background process, with a
set $R\subseteq V$ of \emph{certifiable} states. Each state $q$ has an
authoritative representation $a(q)$, possibly absence $\absent$; a state
$q\in R$ answers a read with a body $b(q)$ and fields $h(q)$, its
\emph{observation} $o(q)=(b(q),h(q))$. Every finite path of $G$ is a
permitted history, and every state may persist for any positive time.
\end{definition}

States outside $R$ stay in $G$. With $R$ the successful current-object reads,
the path $A/g\to\absent\to A/g$ still breaks continuity although its 404 state
supplies no certified observation; conversely $\absent\to A/g\to\absent$
imposes nothing on its endpoints, which are not certifiable. $R=\emptyset$
certifies nothing, and $R=V$ is the all-state special case.

\emph{Read fidelity} requires $b(q)=a(q)$ for every $q\in R$. A cache can
replay an earlier origin response without describing the origin's current
state; such a read fails fidelity whatever its fields, which is a different
failure from an intervening return. If the claim concerns the cache itself,
$a$ must describe the cache.

A \emph{capture} is an atomic read whose instant lies in a known window
$[s,e]$, $s<e$, of one monotone client clock. Captures are \emph{passive}:
they do not change the contract state, so the permitted paths do not depend on
when or whether anyone reads. This is the model of a Web client reading
through public interfaces; it does not claim that a real server keeps no
per-request state, only that no such state constrains the publication paths
the contract permits.

For an observation $\ell$ with body $x$ let $S_\ell=\{q\in R : o(q)=\ell\}$ and
$P_x=\{q\in V : a(q)=x\}$. $\Reach$ and $\CoReach$ are taken in the complete
graph $G$, not in the subgraph induced by $R$, and include paths of length
zero.

\subsection{The continuity criterion}
\begin{theorem}[Continuity criterion]\label{thm:criterion}
Captures returning $\ell$ and later $\ell'$ prove, for every permitted history,
that the representation is $x$ from the first read instant to the second iff
\[\Reach(S_\ell)\cap\CoReach(S_{\ell'})\subseteq P_x\quad\text{in } G.\]
\end{theorem}
\begin{proof}
\emph{Sufficiency.} A permitted history is in some state of $S_\ell$ at the
first read and in some state of $S_{\ell'}$ at the second; every state it
passes in between lies on a path from $S_\ell$ to $S_{\ell'}$, so it is
reachable from $S_\ell$ and co-reachable to $S_{\ell'}$, hence in $P_x$ by the
inclusion. \emph{Necessity.} Let $z$ lie in the intersection with $a(z)\ne x$.
A path from $S_\ell$ to $z$ followed by a path from $z$ to $S_{\ell'}$ is a
finite path, and since every finite path fits into any positive interval with
positive holding times, some permitted history follows it between two read
instants and has representation $a(z)\ne x$ in between. If $z$ is itself an
endpoint, that endpoint read is unfaithful.
\end{proof}

A field is \emph{sound} for a target if the criterion holds for every pair of
equal observations of that target, \emph{reusable} if after every capture the
resource can leave the target and return to the same observation, and
\emph{general} otherwise. Content hashes are reusable under ordinary writes,
and so are S3 version IDs under the documented deletions
(Section~\ref{sec:platforms}).

\subsection{Certificates}
\paragraph{Counterexamples.} For each observation $\ell$ with body $x$, the
product graph has states $(q,d)$ with a sticky departure bit $d$: from
$(q,d)$ an edge $q\to u$ leads to $(u, d\lor a(u)\ne x)$. The initial states
are $(q,\mathit{false})$ for $q\in S_\ell$ and the accepting states
$(q,\mathit{true})$ for $q\in S_{\ell}$. An accepting product path is exactly
a path of $G$ that starts and ends with observation $\ell$ and visits another
representation, so a multi-source breadth-first search returns a shortest
violation for $\ell$, and the minimum over observations is the shortest
overall. With $2|V|$ product states and $2|E|$ product edges per observation,
$k$ observations take $O((k+1)(|V|+|E|))$ time including validation;
fidelity failures, violations of length zero, are checked first. A separate
checker verifies a returned path edge by edge: both endpoints in $R$, equal
endpoint observations, faithful endpoints, and an intermediate state with
another representation. It establishes validity; shortestness follows from
the search and is cross-checked against an independent all-pairs oracle.

\paragraph{Invariants.} When no violation exists, the solver returns for each
observation $\ell=(x,z)$ two sets $B_\ell$ and $D_\ell$ of states reachable
before and after a departure (they may overlap). The checker verifies, without
search: one certificate per observation occurring on $R$; $S_\ell\subseteq
B_\ell$; $a(q)=x$ on $B_\ell$; every edge $q\to u$ with $q\in B_\ell$ leads into
$B_\ell$ if $a(u)=x$ and into $D_\ell$ otherwise; $D_\ell$ is closed under
edges; $D_\ell\cap S_\ell=\emptyset$; and fidelity on $R$. By induction on
path length every path from $S_\ell$ is represented with its correct
departure bit, so no path returns to $S_\ell$ after a departure, self loops
and arbitrary repetitions included. The reachable product states always form
such a certificate, so the format is complete for the finite model.

\paragraph{Bounded transitions.} Let $d$ be the length of a shortest
violation, or $\infty$ if none exists. If a separately justified bound allows
at most $K$ transitions between two reads, continuity holds for every
observation exactly when $K<d$: no violation fits, and the shortest one
witnesses the converse. The bound counts transitions, not elapsed time or
captures. A binary body with a counter modulo~3 needs $d=6$ changes for a
hidden return (equal bodies need an even number, equal counters a multiple
of~3); three bodies that always change give $d=3$; a counter modulo~4 with a
binary body gives $d=4$. A bound derived from counter width alone need
therefore not be tight.

\subsection{Request windows and bridges}
Let two equal certifiable observations of a sound field be read at instants
$u\in[s_1,e_1]$ and $v\in[s_2,e_2]$ with $e_1\le s_2$. Then $u\le e_1\le s_2\le v$, so the resource
keeps its representation on the closed \emph{bridge} $[e_1,s_2]$, the largest
interval guaranteed for every placement of the reads, since $u=e_1$ and $v=s_2$
are possible and the representation may change immediately outside them; for
touching windows it is the single instant $e_1=s_2$. The helper
\code{certified_read_bridge} returns the half-open bridge $[e_1,s_2)$, the
convention of the preregistered replay, and nothing when that is empty, which
means only that no bridge of positive length exists, not that a common instant
is impossible.
The windows must enclose the atomic reads on one monotone client clock;
independently clocked observations need certified error bounds first, and the
HTTP \code{Date} and \code{Last-Modified} headers are not effective times. The
helper \code{common_times} intersects bridges of several resources; an empty
intersection is unresolved rather than negative, because bridges can be
shorter than the true presence intervals.

\subsection{What a certificate does not prove}
A positive certificate proves a property of the supplied graph. It does not
prove that the graph contains every deployment transition, that a cache read
is linearizable to the origin, or that a restart identifier is globally
fresh. If the graph is a conservative abstraction that contains every concrete
trace and preserves observations and representations, a positive result
transfers to the system; a negative abstract path must be concretized before
it shows a real failure. Observing few transitions never justifies removing
permitted ones.

\subsection{Exhaustive checks}
\code{data/constructed/contracts.json} holds 13 constructed contracts that
each exercise one assumption (Table~\ref{tab:contracts}).
\code{scripts/run_contracts.py} computes each verdict, checks every returned
certificate with the separate checker and rejects any disagreement with the
declared verdict. Their mix of safe and unsafe models has no frequency
interpretation; the three fidelity failures are not hidden returns.

\begin{table*}[t]
\caption{The constructed contracts: verdicts and shortest violations, in transitions.}
\label{tab:contracts}
\small
\begin{tabular}{@{}lrrlr@{}}
\toprule
Model & States & Edges & Verdict & Shortest violation \\
\midrule
\texttt{content\_hash\_return} & 2 & 4 & hidden return & 2 \\
\texttt{modulo\_4\_counter} & 8 & 16 & hidden return & 4 \\
\texttt{epoch\_reset\_on\_restart} & 3 & 5 & hidden return & 2 \\
\texttt{boot\_nonce\_reused} & 3 & 5 & hidden return & 2 \\
\texttt{unique\_boot\_nonce} & 3 & 5 & safe & -- \\
\texttt{durable\_epoch} & 3 & 5 & safe & -- \\
\texttt{hash\_with\_no\_return\_policy} & 2 & 3 & safe & -- \\
\texttt{cached\_old\_response} & 3 & 5 & unfaithful read & 0 \\
\texttt{lagging\_replica} & 3 & 5 & unfaithful read & 0 \\
\texttt{non\_atomic\_body\_token} & 3 & 5 & unfaithful read & 0 \\
\texttt{irrelevant\_internal\_transition} & 2 & 4 & safe & -- \\
\texttt{modulo\_3\_binary} & 6 & 12 & hidden return & 6 \\
\texttt{modulo\_3\_ternary} & 9 & 27 & hidden return & 3 \\
\bottomrule
\end{tabular}

\end{table*}

The unit tests compare the implementation with an independent oracle on every
directed graph on three named states with optional self loops ($2^9=512$
graphs), each with every binary body and validator assignment ($8\times8$):
32{,}768 contracts with $R=V$. The oracle computes all-pairs shortest
distances by Floyd--Warshall and searches endpoint pairs with equal
observations and an intermediate state with another body; verdicts and
shortest lengths agree everywhere, and every certificate passes the checker.
A further 5{,}120 contracts enumerate certifiable scopes: every two-state
digraph with binary representations and fields and every $R$ (1{,}024), and
every ordered three-state acyclic graph likewise (4{,}096), with response
bytes outside $R$ deliberately different from the representation. These are
correctness checks of finite instances, not benchmarks.
